%%%%%%%%%%%%%%%%%%%%%%%%%%%%%%%%%%%%%%%%%%%%%%%%%%%%%%%%%%%%%%%%%%%%%%%%%%%%%%%%
%%                           ONDER DE LOEP ARTIKEL                            %%
%%%%%%%%%%%%%%%%%%%%%%%%%%%%%%%%%%%%%%%%%%%%%%%%%%%%%%%%%%%%%%%%%%%%%%%%%%%%%%%%

\artikel{Focus op algebraïsche vaardigheden in de tweede graad}{}{Hilde Eggermont\\ Alexander Holvoet \\ Els Vanlommel}

\startcontents[inhoudloep]	% om inhoud voor het loep artikel te beginnen

\begin{multicols}{2}
	
	\inhoudstafelloep			% om inhoudstafel voor het loep artikel te tonen
	
	%% Gebruik \section{titel} en \subsection{titel} om genummerde tussentitels te maken
\section{Inleiding}
\subsection{De aanleiding: nieuwe minimumdoelen wiskunde}
Sinds de invoering (vanaf 2019) van de nieuwe eindtermen/minimumdoelen voor wiskunde komen algebraïsche rekenvaardigheden in het curriculum minder expliciet aan bod dan voordien het geval was. Er wordt voor het rekenwerk meer ingezet op het gebruik van ICT. Leraren merken in de praktijk echter dat het vlot kunnen inzetten van algebraïsche vaardigheden erg belangrijk is, ook voor het begrijpen en oplossen van complexere wiskundige problemen. Het oplossen van algebraïsche problemen vereist nauwkeurigheid, patroonherkenning en analytisch denken. Door te oefenen op algebraïsche vaardigheden train je leerlingen dus ook op deze andere vlakken. Uiteraard weet elke wiskundige dat het rekenwerk niet het hoofddoel is van wiskunde en zelden gaat een wiskundehart er sneller van kloppen, maar leerlingen die onvoldoende vaardig zijn in algebraïsch rekenen, lopen vaak vast bij wat wél essentieel is, zoals het modelleren van problemen, het oplossen van vraagstukken, het begrijpen van nieuwe wiskundige concepten\ldots

\subsection{Opbouw van deze loep}
Na deze inleiding, kaderen we de loep in paragraaf 2 eerst in een groter geheel. Je vindt in dit deel verwijzingen naar extra achtergrondliteratuur. In paragraaf 3 gaan we na hoe vakgroepen kunnen werken aan een leerlijn voor algebraïsche vaardigheden en in de volgende twee paragrafen maken we dat concreet door te kijken naar ontbinden in factoren (paragraaf 4) en rekenen met rationale lettervormen (paragraaf 5), wat vroeger in de tweede graad behandeld werd. In paragraaf 6 focussen we op het gericht inzetten van ICT voor remediëring, zowel voor de tweede als de derde graad. We lijsten hierin de verschillende functionaliteiten op en bespreken ICT-tools die we kunnen inzetten bij concrete vragen van leerlingen.


\subsection{Voor wie is deze loep geschikt?}
We richten ons in deze loep tot vakgroepen die willen werken aan een leerlijn rond algebraïsche vaardigheden in hun school. Die vaardigheden zijn voor alle leerlingen belangrijk, al hangen de noden natuurlijk af van het aantal uren wiskunde in de studierichting. Het is een taak van vakgroepen en leraren om hierin doordachte keuzes te maken. Met deze loep willen we inspiratie bieden bij het maken van die keuzes.


\section{Ruimer kader}

Om te beginnen, willen we sterk benadrukken dat de focus binnen wiskunde uiteraard niet louter op procedureel algebraïsch rekenen mag liggen en dat we met deze loep geenszins pleiten voor een reductie tot receptenwiskunde. Procedureel en conceptueel wiskundeonderwijs zijn immers onlosmakelijk met elkaar verweven. We verwijzen in dit kader naar de vijf strengen van wiskundige bekwaamheid van Kilpatrick et al. (2001), zoals afgebeeld in figuur \ref{fig:strengen}.

\afbeelding[0.8\linewidth]{img/LAlgebra/LEls/strengen_Kilpatrick}{De vijf strengen van wiskundige bekwaamheid}{fig:strengen}

%\url{https://www.fi.uu.nl/wiskrant/artikelen/314/314juni_zwaneveld.pdf}

Automatiseren is meer dan alleen maar meer oefenen. Belangrijk is na te denken over hoe je algebraïsche vaardigheden didactisch best aanbrengt. We gaan daar in deze loep niet verder op in en verwijzen naar de loep van Uitwiskeling 24/1 (winter 2008) over algebra met applets. Daarin hebben we al uitgebreid met voorbeelden geïllustreerd hoe je bij het aanbrengen van de leerstof inzicht in algebra kunt ontwikkelen. Het ging in die loep bijvoorbeeld over het introduceren van letters om te veralgemenen, meetkundig voorstellen van bewerkingen en leren oplossen van tweedegraadsvergelijkingen. Het is zeker de moeite waard om deze loep terug ter hand te nemen.

Verder mogen we natuurlijk niet uitsluitend inzetten op veelvuldig oefenen, maar moeten
we integendeel een goede combinatie nastreven van rekenvaardigheid en inzicht in wat er moet gebeuren. We kennen zelf ook wel uit onze eigen ervaring de voorbeelden van leerlingen die een vergelijking als $-2(x-7)=(x-7)(x+3)$ oplossen door eerst de haakjes uit te werken en dan opnieuw te ontbinden in factoren. Of leerlingen die bij het zien van een kwadratische uitdrukking in een pavloviaanse reactie meteen een discriminant berekenen, ook als die helemaal niet nodig is. Of nog: leerlingen die machteloos staan als ze vaststellen dat ze een formule vergeten zijn. Om zulke zaken te vermijden, is het belangrijk ervoor te zorgen dat je bij het oefenen een beroep blijft doen op inzicht. Ook dat aspect behandelen we niet in deze loep. We verwijzen graag naar de felgesmaakte loep van Uitwiskeling 29/1 (winter 2013) rond algebra oefenen met inzicht waarin we veel kleine voorbeelden geven van verstandige manieren om algebra te oefenen.


Wat ervaren leraren weten uit hun lespraktijk wordt bevestigd door allerhande studies die het belang benadrukken van geautomatiseerde basisvaardigheden zoals algebraïsch rekenen. We sommen hier enkele verwijzingen op naar artikels over dit thema. Ze kunnen dienen als bijkomende inspiratie om een eigen visie te ontwikkelen.
\begin{itemize}
	\item Net zoals beginnende sporters en muzikanten talloze uren spenderen aan het automatiseren van 	basishandelingen, is het belangrijk dat leerlingen voldoende tijd kunnen investeren in het verwerven van algebraïsche rekenvaardigheden. Een matige beheersing volstaat niet, zoals goed wordt verwoord in het rapport van een adviesgroep voor het verbeteren van het wiskundeonderwijs in de Verenigde Staten (Geary et al., 2008). De auteurs spreken in dit rapport van ‘over-leren’: \emph{`[I]n support of complex problem solving, arithmetic facts and fundamental algorithms should be thoroughly mastered, and indeed, over-learned, rather than merely learned to a moderate degree of proficiency.'}
	\item Bergsen (2023) beschrijft hoe leerlingen op drie basisscholen in zes weken beter leerden rekenen door dagelijks rekenhandelingen te automatiseren en rekenfeiten te memoriseren. In het artikel staan heel wat praktische tips en verwijzingen naar onderzoek rond het belang van automatiseren die ook relevant zijn voor het secundair onderwijs.
	\item Sinds 1995 hebben alle staten in de VS standaarden aangenomen voor het wiskundeonderwijs (basis en secundair). In de meeste staten is de nadruk hierbij minder gelegd op het memoriseren en meer op redeneren op basis van concepten. Deze verandering ging ervan uit dat wiskundig redeneren mogelijk is zonder terug te vallen op kennis opgeslagen in het geheugen. De resultaten bleven echter uit en cognitieve wetenschappers hebben ondertussen vastgesteld dat deze veronderstelling onjuist was. Ons werkgeheugen is heel beperkt. Daarom zijn goed gememoriseerde feiten en geautomatiseerde routines nodig bij het oplossen van wiskundige problemen, ongeacht de complexiteit ervan. Hartman et al. (2023) doen wetenschappelijk onderbouwde voorstellen om de wiskundeprestaties van leerlingen te verbeteren.	
	\item De relevantie van algebra en algebraïsche vaardigheden in het curriculum wordt krachtig samengevat door Tall en Thomas (1991) in een artikel waarin het gaat over het aanleren van algebraïsche vaardigheden met de computer: \emph{`There is a stage in the curriculum when the introduction of algebra may make simple things hard, but not teaching algebra will soon render it impossible to make hard things simple.'}   De auteurs benadrukken het belang van `flexibel' denken in algebra, een soort helikopteroverzicht bewaren om bij elk probleem de gepaste rekenstrategie toe te kunnen passen en dus te rekenen met inzicht.
\end{itemize}





\section{Vragenlijst voor de vakgroep}

Wanneer je extra wilt inzetten op algebraïsche rekenvaardigheden, moet je met de vakgroep een aantal keuzes maken. In deze paragraaf gaan we hier dieper op in.

Het komt erop neer dat je met de vakgroep nadenkt over een soort van verticale leerlijn rond rekenvaardigheden. Hierbij wil je over de leerjaren en graden heen een consistente competentieopbouw garanderen. Dit zorgt voor houvast voor elke vakleerkracht: je weet wat de beginsituatie is van je leerlingen, je weet waar je naartoe moet, wat de samenhang is tussen de doelen en de opbouw van de doelen. Bovendien zorg je er met een leerlijn voor dat voldoende herhaling aan bod komt, maar tegelijkertijd ook dat overbodige herhalingen en overlappingen worden vermeden.

Om je te helpen bij het opstellen van zo'n leerlijn hebben we een aantal concrete vragen opgelijst. Dit is geen stappenplan dat je van a tot z moet doorlopen, maar dient louter ter inspiratie. Het is ook niet de bedoeling dat dit voor een berg extra planlast zorgt. Belangrijk is dat je tot een gedeelde visie en aanpak komt.

\end{multicols}

\begin{kader}{}{Inspiratievragen voor de vakgroep}
	\begin{itemize}[label=\textcolor{\rubriekkleur}{\Large\textbullet}]
	\item Waarom willen we inzetten op rekenvaardigheden? Wat is het uiteindelijke doel? In welke mate vinden we rekenvaardigheden belangrijk? Hebben we hierover een gemeenschappelijke visie met de vakgroep?
	\item Welke algebraïsche rekenvaardigheden en technieken zijn belangrijk in de richtingen waarvoor de leerlijn wordt opgesteld? Je kunt hierbij terugvallen op de doelen en leerinhouden uit het leerplan dat je gebruikt.
	\item Met welk doel behandelen we een bepaalde rekentechniek? Is dat echt nodig? Hoe ver gaan we daarin? Kunnen we daarin ook te ver gaan? Wat zijn de valkuilen? Kunnen we iets schrappen?
	\item Zien we een lijn, opbouw en samenhang tussen de opgelijste inhouden? Visualiseer dit eventueel in een schema.
	\item Over welke jaren en graden bouwen we de leerlijn op? Wat komt concreet in elk jaar aan bod? Op welke plaats in het jaarplan bouwen we dit in?
	\item Hoeveel tijd voorzien we voor het aanbrengen en inoefenen van de vaardigheden?
	\item Hoe sterk moeten de rekenvaardigheden geautomatiseerd worden? Hoe doen we dat concreet?
	\item Bouwen we doorheen de leerlijn voldoende herhaling in zodat we de vaardigheden consolideren? Dezelfde vaardigheden kunnen bij verschillende leerinhouden terugkomen.
	\item Welke vaardigheden en inhouden zijn essentieel en welke maken eventueel deel uit van een verdiepingspakket?
	\item Hoe kunnen we werken aan remediëring als het niet loopt zoals gepland? Welke ondersteuning kunnen we aan leerlingen bieden?
	\item Hoe evalueren we de inhouden en vaardigheden? Hoe gaan we na of we de vooropgestelde doelen bereikt hebben?
	\item Hoe verduidelijken we aan leerlingen wat we verwachten?
	\item Welk bestaand lesmateriaal kunnen we gebruiken om de doelen te bereiken? Moeten we extra lesmateriaal ontwikkelen?
	\end{itemize}
\end{kader}


\begin{multicols}{2}

In de volgende twee paragrafen bekijken we concreet waar we ontbinden in factoren en rekenen met rationale lettervormen in het curriculum aan bod kunnen laten komen. Omdat deze leerstofonderdelen niet meer in alle handboekenreeksen aan bod komen, voegden we telkens een aantal mogelijke voorbeeldoefeningen toe.

 
\section{Ontbinden in factoren}

Eén van de meest gebruikte technieken bij het algebraïsch rekenwerk is misschien wel het ‘ontbinden in factoren’. Je hebt het bijvoorbeeld nodig als je breuken of wortels wilt vereenvoudigen. Ook het aflezen van nulwaarden is een stuk eenvoudiger als de uitdrukking ontbonden is in factoren. 
\end{multicols}
 \afbeelding[0.9\textwidth]{img/ontbinden.jpg}{}{}
\begin{multicols}{2}
Er is een hele resem van technieken die je hiervoor kunt aanwenden. Het is logisch dat de kennis en de oplossingsmethoden waarover de leerlingen beschikken groeit in de loop van hun loopbaan. Het lijkt ons geen goed idee om te wachten tot alle technieken voorhanden zijn om te beginnen oefenen op een basisvaardigheid als ontbinden in factoren. Een spiraalaanpak is hier zeker op zijn plaats. Vandaar een warm pleidooi om er reeds in het derde jaar mee te starten. De technieken die de leerlingen dan voorhanden hebben zijn het ‘buiten de haakjes brengen’ (omgekeerde distributiviteit) en enkele merkwaardige producten. Dit is kennis die de leerlingen in de eerste graad hebben opgedaan, maar het accent ligt dan vooral op het ‘uitwerken’ ( $(a+b)(a-b)=a^2-b^2$ ). De andere richting (ontbinden in factoren) komt nagenoeg niet aan bod. 
	\subsection{Wat al kan in de eerste graad}

Hoewel de minimumdoelen in de eerste graad alleen het uitwerken vragen en niet het omgekeerde, namelijk het ontbinden in factoren, lijkt het ons goed om dit toch aan bod te laten komen. Dit kan in de les en hoeft niet noodzakelijk getoetst te worden bij grote evaluaties. In de lesactiviteit hieronder geven we opdrachten die in dat kader gemaakt kunnen worden. We haalden onze inspiratie uit Gevers, P., e.a. (1970). In de eerste opgave wordt er ook al met letters gewerkt. De tweede opdracht is nauw verwant met het efficiënt rekenen van de lagere school. Op die manier laten we de vaardigheden uit de lagere school niet verloren gaan en leggen we de connectie met de nieuwe leerstof. Het rekenen met negatieve getallen komt er als nieuwe moeilijkheid bij.  Afhankelijk van de opbouw van de leerstof zijn dit oefeningen voor het eerste of het tweede jaar. 


\end{multicols}
\begin{lesactiviteit}{Sommen en verschillen omzetten in producten}

\begin{vraagenantwoord}
\vraag{Gebruik de distributieve eigenschap om elk van de volgende sommen (of verschillen) te schrijven als een product.}

%\begin{vraagenantwoord}
\begin{itemize}[label=\textcolor{\rubriekkleur}{\Large\textbullet}]

\begin{multicols}{2}
    
    \item {$2\cdot 3 + 2\cdot 4=$}
    \item {$7\cdot 2 - 7\cdot 9=$}
    \item {$-18 + 6\cdot (-9)=$}
    \item {$2x+4=$}
    \item {$-4c + 20b=$}
    \item {$ab + (-2)b=$}
    \item {$(-a)(-b) + 2ab=$}
    \item {$abc + bc=$}
\end{multicols}
\end{itemize}
%\end{vraagenantwoord}


\vraag{Schrijf de volgende sommen (of verschillen) op als producten en reken daarna uit zonder rekentoestel.}
%\begin{vraagenantwoord}
\begin{itemize}[label=\textcolor{\rubriekkleur}{\Large\textbullet}]   
\begin{multicols}{2}
\item {$14\cdot 17 + 16\cdot 17 =$}
\item {$7\cdot 11 - 7\cdot 19 =$}
\item {$38\cdot (-12) + (-7)\cdot (-38) =$}
\item {$72\cdot 13 + (-36)\cdot 16 =$}
\item {$5\cdot (-3) - (-3)\cdot (-7) + 8\cdot (-3) =$}
\item {$(-2)\cdot (-7) + (-1)\cdot (-7) - 8\cdot 7 =$}
\end{multicols}
%\end{vraagenantwoord}
\end{itemize}

\end{vraagenantwoord}
\end{lesactiviteit}
\begin{multicols}{2}


\subsection{Uitbouwen technieken in de tweede graad}
\subsection*{In het derde jaar}
\label{q:efficientie}
In het derde jaar komen er op zich weinig nieuwe technieken bij. De opgaven worden een beetje complexer en de leerlingen moeten vooral oefenen. Er wordt meer met letters gerekend en de formules van merkwaardige producten kunnen aangevuld worden met de formules voor een derde macht van een som of verschil. Deze formules staan niet in de minimumdoelen vermeld, maar een leerling uit een doorstroomrichting zou minstens in staat moeten zijn om een uitdrukking als $(a+b)^3$ correct uit te werken, al dan niet met de formule. De omgekeerde weg, namelijk het ontbinden van $a^3+3a^2b+3ab^2+b^3$ in factoren, is lastiger en vraagt wel wat oefening. Hiervoor is het wel nodig om de formule te kennen. 

We starten de lessenreeks met een opgave waar de nadruk ligt op het herkennen van sommen en producten. Door stil te staan bij deze begrippen kunnen we de leerlingen duidelijk maken wat we bedoelen met ´ontbinden in factoren' of ´schrijf als een product'. Zo krijgen ze inzicht in de opbouw van een formule en de `grammatica' van de algebra. Het gebruik van de correcte terminologie helpt hierbij. In de lesactiviteit hieronder vind je twee opgaven die daarop focussen. Het is niet de bedoeling om hier ellenlange lijsten van oefeningen op te maken, maar het is wel zinvol om doorheen je lessen dit geregeld aan te halen. Zo kun je bijvoorbeeld bij de stelling van Pythagoras even stilstaan bij de vorm van de twee leden van de formule: het ene is een macht, het andere een som (van machten). 

Inzicht in de opbouw van wiskundige uitdrukkingen is belangrijk. Als je breuken moet optellen bijvoorbeeld, is het handig als de noemers ontbonden zijn in factoren. Het vinden van het kleinste gemeenschappelijke veelvoud is dan immers een stuk eenvoudiger. Het herkennen van de structuur van een wiskundige uitdrukking komt ook later nog van pas, bijvoorbeeld om bij het afleiden de juiste formule te gebruiken. Inzicht moet je opbouwen en inoefenen. Vandaar dat we er van in het derde jaar aan werken. De lesactiviteit is gebaseerd op de bundel die het team van het Wiskundeplan ter beschikking stelt op haar website \url{www.wiskundeplan.be}. 


\end{multicols}

 \afbeelding{img/zonnecreme.jpg}{}{}

\begin{lesactiviteit}{Sommen en producten}
	
	\begin{vraagenantwoord}
		\vraag{Is de gegeven uitdrukking te beschouwen als een product of als een som?}
  \begin{multicols}{2}
    %\begin{vraagenantwoord}	
    \begin{itemize}[label=\textcolor{\rubriekkleur}{\Large\textbullet}]
		\item {$a\cdot b\cdot c +1$}
        \item {$(x+y)^2$}
        \item {$a\cdot (4b+c)$}
        \item {$x\cdot y + x\cdot z$}
        \item {$a\cdot b\cdot (c +1)$}
        \item {$(x+y)(x-z)$}
	\end{itemize}	
      
  \end{multicols}
     % \end{vraagenantwoord}
  \begin{multicols}{2}
  \begin{vraagenantwoord}
       \antwoord{som}
       \antwoord{product}
       \antwoord{product}
       \antwoord{som}
       \antwoord{product}
       \antwoord{product}
  \end{vraagenantwoord}
        
  \end{multicols}
          
  \vraag{Onderzoek telkens of je van links naar rechts ontbindt, uitwerkt of geen van beide.}
     \begin{multicols}{2}
         
     
     \begin{itemize}[label=\textcolor{\rubriekkleur}{\Large\textbullet}]
		\item {$(3a+b)^2=9a^2+6ab+b^2$}
        \item {$4(x^2 -y^2)=4x^2-4y^2$}   
        \item {$4(x^2-y^2)=4(x-y)(x+y)$}
        \item $3x^3 -3x+1 =3x(x^2-1)+1 \\=3x(x-1)(x+1)+1$
        \item {$27a^3-27a^2b+9ab^2-b^3=(3a-b)^3$}                 
        \item {$4ax +2z+2az+2x\\=2x(2a+1)+2z(1+a)$}
     \end{itemize}

   \end{multicols}
   \begin{multicols}{2}
      \begin{vraagenantwoord}
          \antwoord{uitwerken}
          \antwoord{uitwerken}  
          \antwoord{ontbinden}
          \antwoord{geen van beide}
          \antwoord{ontbinden}
          \antwoord{geen van beide}
      \end{vraagenantwoord}
  \end{multicols} 
  
\end{vraagenantwoord}       
\end{lesactiviteit}


\begin{multicols}{2}
    
In de volgende lesactiviteit geven we eerst een overzicht van de formules, met telkens een voorbeeld erbij. In tegenstelling tot de eerste graad ligt nu de nadruk niet op het uitwerken, maar op het ontbinden in factoren. Aangezien dit nog niet is ingeoefend in de eerste graad, vraagt dit heel wat oefening om het in de vingers te krijgen. Je laat de leerlingen best enkele opgaven in de klas maken zo dat je onmiddellijk feedback kunt geven, maar het zal ook nodig zijn om thuis hier verder op te oefenen. Voorzie dus voldoende materiaal. In paragraaf 6 geven we heel wat tips om hiervoor ICT in te schakelen. 

De eerste lesactiviteit is gericht op het inoefenen van de aparte technieken. De technieken zelf zijn in vorige lessen aangebracht en staan samengevat bij de start van de lesactiviteit. Inspiratie voor de aanbreng kun je bijvoorbeeld vinden in de loep van Uitwiskeling 29/1 (winter 2013). In de tweede lesactiviteit komen dan gemengde oefeningen aan bod. 
\end{multicols}
\clearpage


\begin{lesactiviteit}{Technieken om te ontbinden in factoren}

\begin{kader}{}{Overzicht}
		Ontbinden in factoren betekent: een som (van termen) omzetten in een product (van factoren). We kunnen daarvoor gebruik maken van de volgende technieken. Een voorbeeld staat telkens in de rechterkolom.
  \begin{tabel}
	\begin{tabular}{l l|l}
    \rowcolor{\rubriekkleur!10} 1. &  Afzonderen & $24x^3-15x^2+3x=3x(8x^2-5x+1)$\\
    \hline
    \rowcolor{\rubriekkleur!10} 2. &  Merkwaardige producten: & \\
    \rowcolor{\rubriekkleur!10} & $a^2-b^2=(a-b)(a+b)$ & $25x^2-3=(5x-\sqrt{3})(5x+\sqrt{3})$\\
    \rowcolor{\rubriekkleur!10} & $a^2\pm 2ab + b^2 = (a\pm b)^2$ & $4x^2+12xy+9y^2=(2x+3y)^2$\\
    \rowcolor{\rubriekkleur!10} & $a^3 \pm 3a^2b +3ab^2 \pm b^3 = (a \pm b)^3$ & $t^3-6t^2+12t-8=(t-2)^3$\\
    \hline
    \rowcolor{\rubriekkleur!10} 3. & Twee aan twee samennemen & $3x^3-6x^2 +x-2=(3x^3-6x^2)+(x-2)$\\
    \rowcolor{\rubriekkleur!10} & & \phantom{$3x^3-6x^2 +x-2$} $=3x^2(x-2)+(x-2)$\\
    \rowcolor{\rubriekkleur!10} & & \phantom{$3x^3-6x^2 +x-2$} $=(x-2)(3x^2+1)$\\
 \end{tabular}
 \end{tabel}
        Afspraak: we ontbinden altijd zo ver mogelijk in factoren van de eerste en/of tweede graad.
\end{kader}

\begin{vraagenantwoord}
    \vraag{Ontbind in factoren door af te zonderen.}
    \begin{multicols}{2}
        
    %\begin{vraagenantwoord}
    \begin{itemize}[label=\textcolor{\rubriekkleur}{\Large\textbullet}]
        \item {$\sqrt{5} a-3 \sqrt{5} b$}
        \item {$\sqrt{7}x^2+2\sqrt{7}$}
        \item {$15a^2b+12ab-9ab^3$}
        \item {$x(x+4)-2(x+4)$}
        \item {$2a(3a+1)+5b(3a+1)$}
        \item {$(y^2+3)(y-2)+y-2$}
        \item {$(a+4)(3a+2)-3a-2$}
        \item {$(x-3)(2x+3)-(4x+1)(2x+3)$}
        \item {$3(x-1)+(x-1)^2$}
        \item {$5(a+b)m-4(a+b)n$}
    \end{itemize}
    %\end{vraagenantwoord}

    \end{multicols}
    \vraag{Vul aan tot een merkwaardig product, indien mogelijk, en ontbind in dat geval.}
    \begin{multicols}{2}
        
    %\begin{vraagenantwoord}
    \begin{itemize}[label=\textcolor{\rubriekkleur}{\Large\textbullet}]
        \item {$9x^2+\ldots x+16$}
        \item {$a-6ab+\ldots$}
        \item {$\ldots y^4-42y^2+9$}
        \item {$121x^2-\ldots x^3+4$}
        \item {$-81-4x^2+ \ldots $}
        \item {$x^4+ \ldots + \frac{y^2}{4}$}
      \end{itemize}  
    %\end{vraagenantwoord}

    \end{multicols}
    \vraag{Schrijf de tweetermen als een product indien mogelijk. Indien het niet mogelijk is, leg uit waarom niet.}
    \begin{multicols}{2}
    \begin{itemize}[label=\textcolor{\rubriekkleur}{\Large\textbullet}]     
   % \begin{vraagenantwoord}
        \vraag{$x^2-121$}
        \vraag{$9x^2+16$}
        \vraag{$16x^2-25$}
        \vraag{$36a^2-5$}
        \vraag{$4a^2-81b^2 $}
        \vraag{$144x^2-49y^4$}
        \vraag{$5x^2-3$}
        \vraag{$-4x^2+25y^2$}
     
  %  \end{vraagenantwoord}
    \end{itemize}
    \end{multicols}
\clearpage
    \vraag{Ontbind de drietermen in factoren.}
    \begin{multicols}{2}
    \begin{itemize}[label=\textcolor{\rubriekkleur}{\Large\textbullet}]        
   % \begin{vraagenantwoord}
        \vraag{$25x^2-20x+4$}
        \vraag{$x^4+6x^2+9$}
        \vraag{$25a^2-30ab+9b^2$}
        \vraag{$81+126t+49t^2$}
      
    %\end{vraagenantwoord}
    \end{itemize}
    \end{multicols}

\vraag{Ontbind de volgende viertermen in factoren.}
    \begin{multicols}{2}
    \begin{itemize}[label=\textcolor{\rubriekkleur}{\Large\textbullet}]        
   % \begin{vraagenantwoord}
        \vraag{$a^3+6a^2b+12ab^2+8b^3$}
        \vraag{$8x^3-36x^2y+54xy^2-27y^3$}
        \vraag{$1000x^3+1200x^2b+480xb^2+64b^3$}
        \vraag{$27a^3x^6-54a^2bx^4+36ab^2x^2-8b^3$}
      
    %\end{vraagenantwoord}
    \end{itemize}
    \end{multicols}
     \vraag{Ontbind de viertermen in factoren door termen samen te nemen.}
   
     \begin{itemize}[label=\textcolor{\rubriekkleur}{\Large\textbullet}]        
   % \begin{vraagenantwoord}
     \begin{multicols}{2}
     
        \vraag{$4x+4y+ax+ay$}
        \vraag{$3x+5y+6xy+10y^2$}
        \vraag{$2ax+3a-2bx-3b$}
        \vraag{$7ab+14a+b+2$}
        \vraag{$1-a-b+ab$}
        \vraag{$5x^3-x^2+35x-7$}
        \vraag{$x^3+x^2+x+1$}
        \vraag{$2x^6+3x^4+6x^2+9$}
        \vraag{$a^2-6b-2ab+3a$}
        \vraag{$12ab^2-1-4b^2+3a$}
     \end{multicols}
    %\end{vraagenantwoord}
     \end{itemize}
   
      \vraag{Met welke techniek begin je de ontbinding in factoren bij de volgende uitdrukkingen? Je moet de ontbinding (nog) niet helemaal uitwerken.}
    \begin{multicols}{2}
    \begin{itemize}[label=\textcolor{\rubriekkleur}{\Large\textbullet}]       
   % \begin{vraagenantwoord}
       \vraag{$(x-1)(x+3)-5(x+3)$}
       \vraag{$-1+7x^2$}
       \vraag{$49-14x^2+x^4$}
       \vraag{$(x-z)y-z(z-x)$}
       \vraag{$4x^4-9y^2$}
       \vraag{$4(a^2-b)^2-9(a^2-3b)^2$}
       \vraag{$x^4+\frac{1}{4}+x^2$}
       \vraag{$12ab-8a+3b-2$}
   %    \end{vraagenantwoord}
   \end{itemize}
    \end{multicols}
    \begin{multicols}{2}
       \begin{vraagenantwoord}
           \antwoord{Factor $(x+3)$ afzonderen.}
           \antwoord{Verschil van twee kwadraten.}
           \antwoord{Merkwaardig product: kwadraat van een verschil.}
           \antwoord{Factor $(x-z)$ afzonderen.}
           \antwoord{Verschil van twee kwadraten.}
           \antwoord{Verschil van twee kwadraten.}
           \antwoord{Merkwaardig product: kwadraat van een som.}
           \antwoord{Termen twee aan twee samennemen: bv. de eerste twee en de laatste twee. }
       \end{vraagenantwoord}    
    \end{multicols}
\end{vraagenantwoord}
    
\end{lesactiviteit}
 \afbeelding[0.8\textwidth]{img/ontbindingslus.jpg}{}{}
\begin{lesactiviteit}{Ontbinden in factoren: gemengde oefeningen}
\begin{vraagenantwoord}
    
     \vraag{Ontbind zo ver mogelijk in factoren. Dit wil zeggen: schrijf de uitdrukkingen als een product van factoren van de eerste en/of tweede graad die je niet verder kunt ontbinden.}
    \begin{multicols}{2}
        
    %\begin{vraagenantwoord}
     \begin{itemize}[label=\textcolor{\rubriekkleur}{\Large\textbullet}] 
        \vraag{$4a^3-25a$}
        \vraag{$x^4-1$}
        \vraag{$18y^5-22y^3$}
        \vraag{$32a^3b-50ab^3$}
        \vraag{$144x^3-x$}
        \vraag{$16a^4-9b^4$}
        \vraag{$-a^2-2ab-b^2$}
        \vraag{$4x^5-28x^4+49x^3$}
        \vraag{$8x^4y-24x^3y^3+18x^2y^5$}
        \vraag{$(2x^2+3)(x+1)-(x^2+4)(x+1)$}
      \end{itemize}  
    %\end{vraagenantwoord}

    \end{multicols}

    \vraag{Ontbind in factoren.}
    \begin{multicols}{2}
        
   % \begin{vraagenantwoord}
    \begin{itemize}[label=\textcolor{\rubriekkleur}{\Large\textbullet}] 
        \vraag{$(x+2y)^2-16$}
        \vraag{$(x+1)^2-6(x+1)+9$}
        \vraag{$4(2x+3)^2-20(2x+3)+25$}
        \vraag{$9(5x-y)^2-16(3x+2y)^2$}
     \end{itemize} 
    %\end{vraagenantwoord}

    \end{multicols}

    \vraag{Ontbind in factoren. Soms kun je op verschillende manieren starten. Probeer dan een doordachte keuze te maken.}
    \begin{multicols}{2}
     \begin{itemize}[label=\textcolor{\rubriekkleur}{\Large\textbullet}]     
    %\begin{vraagenantwoord}
        \vraag{$9x^2-81x^2y^2$}
        \vraag{$x^3+x^2-4x-4$}
        \vraag{$-x^6-4x^4-4x^2$}
        \vraag{$a^2-(a+b+c)^2$}
        \vraag{$x^8-81$}
        \vraag{$a^2(a^2-3)-4(a^2-3)$}
        \vraag{$(2x-2)(3x+4)-(x-1)(2x-1)$}
        \vraag{$-\sqrt{14}x^3y+\sqrt{21}x^2y+\sqrt{35}x^2y^2$}
        \vraag{$-18s^2t^3+84s^2t^2-98s^2t$}
        \vraag{$ac^2+bd^2-ad^2-bc^2$}
        \vraag{$(3x+4)^2-8(3x+4)+16$}
        \vraag{$x^3-x^2-x+1$}
        \vraag{$2x^3-13x^2-6x+39$}
        \vraag{$(5x+1)^2+4(5x+1)$}
        \vraag{$ax^2+bx^2-4a-4b$}
        \vraag{$(4x^2+1)^2-(12x-8)(4x^2+1)$}
        \vraag{$\sqrt{8}-ab+\sqrt{2}b-2a$}
        \vraag{$a^3+b^3+a^2b^2+ab$}
        \vraag{$(a+2b)^2-(2a-b)^2$}
        \vraag{$(a-b)^3+(a+b)^3$}
     \end{itemize}   
   % \end{vraagenantwoord}

    \end{multicols}
\vraag{Hieronder zie je twee manieren om de voorlaatste opgave van de vorige oefening aan te pakken. De resultaten zijn verschillend. Geven beide werkwijzen een correct antwoord? Leg uit.}
Werkwijze 1
\begin{align*}
    (a+2b)^2-(2a-b)^2 &= (a+2b-(2a-b))(a+2b+2a-b)
    \\&= (-a+3b)(3a+b)
\end{align*}
Werkwijze 2
\begin{align*}
    (a+2b)^2-(2a-b)^2 &= a^2+4ab+4b^2-(4a^2-4ab+b^2)
    \\&= -3a^2+8ab+3b^2
\end{align*}
\antwoord{De enige goede manier van werken is de eerste. Je herkent in de opgave een verschil van kwadraten. Door de formule van dat merkwaardig product te gebruiken ontbind je in factoren zoals gevraagd en kom je tot de oplossing. Bij de tweede werkwijze worden eerst de kwadraten uitgewerkt. Dit geeft helemaal geen antwoord op de vraag omdat het resultaat een som is en geen product. In principe kun je die drieterm ontbinden, maar daar zijn technieken uit het vierde jaar voor nodig.}

\end{vraagenantwoord}
\end{lesactiviteit}

\begin{multicols}{2}
   
De beheersing van de technieken van het ontbinden in factoren is niet het einddoel van deze lessenreeks. De bedoeling is om het later te kunnen inzetten bij bijvoorbeeld het rekenen met letterbreuken of het rekenen met goniometrische uitdrukkingen. Om de technieken in de vingers te houden is het aangeraden om regelmatig eens een kleine opdracht in te lassen. Anders dreigen de vaarigheden verloren te gaan en dan is de tijd die we er in gestopt hebben verloren tijd, wat erg jammer zou zijn.


\subsection{In het vierde jaar}

In het vierde jaar breiden we de set aan technieken verder uit: ontbinden van een algemene drieterm van de tweede graad en ontbinden door gebruik te maken van de reststelling. De lijst van merkwaardige producten wordt aangevuld met de formules voor $a^3 \pm b^3$. Deze leerstof stond voor 2022 ook op het leerplan van het vierde jaar en is dus terug te vinden in de handboeken van toen. Sommige handboeken hebben dit al terug opgenomen in hun aangepaste uitgave, andere werken met extra katernen. We gaan hier niet in op het aanbrengen van deze leerstof. In paragraaf 6 vind je suggesties om via ICT materiaal aan te bieden aan de leerlingen om al dan niet zelfstandig verder te oefenen. 

Het is niet zo duidelijk wat het leerplan van het Katholiek Onderwijs Vlaanderen precies verwacht op het vlak van het ontbinden van tweedegraadsveeltermen. Het wordt vermeld bij de wenken in het kader van het oplossen van tweedegraadsvergelijkingen. Het lijkt een gemiste kans om hier toch niet iets verder te gaan en niet alleen onvolledige tweedegraadsvergelijkingen op te lossen door het ontbinden in factoren. Maar ook omgekeerd, het oplossen van tweedegraadsvergelijkingen in te zetten om een algemene tweedegraadsuitdrukking te ontbinden in factoren. 

In richtingen met vijf uur wiskunde komen in het vierde jaar nog andere technieken aan bod zoals de regel van Horner in het hoofdstuk over veeltermen. In de meeste handboeken is hiervoor voldoende materiaal ter beschikking.

	
\section{Rekenen met rationale lettervormen}


Naast ontbinden in factoren komt ook het rekenen met breuken en letters vaak aan bod. Denk maar aan het vereenvoudigen van functievoorschriften waarin een rationale vorm voorkomt, het rekenen met machten, bewijzen van goniometrische identiteiten, berekenen van integralen van rationale functies, splitsen in partieelbreuken\ldots Dit is leerstof die vooral in de derde graad voorkomt, voornamelijk in richtingen met meer uren wiskunde. 

In deze paragraaf gaan we dieper in op het rekenen met rationale lettervormen en het concrete inoefenen daarvan in het vierde jaar.




\subsection{Schets van de verticale opbouw}
Ook voor het rekenen met rationale lettervormen denken we dat het goed is om met kleine stappen te werken en een coherente verticale opbouw op te stellen die aanleiding geeft tot voldoende veel herhaling en inoefening. Het is voor leerlingen veel moeilijker wanneer ze in de derde graad tegelijk nieuwe concepten én het algebraïsch rekenwerk moeten aanleren. Om hen optimaal kansen te geven, zetten we al vroeger in op de nodige algebraïsche vaardigheden.

In de eerste graad komt het rekenen met breuken en het rekenen met letters aan bod. Ontbinden in factoren wordt, zoals in de vorige paragraaf werd besproken, in het tweede jaar voorbereid door het daar al een eerste keer aan te halen bij de leerstof van merkwaardige producten. Dat ontbinden wordt expliciet verder ingeoefend in het derde jaar en komt terug in het vierde jaar bij tweedegraadsvergelijkingen (of -functies). In richtingen met vijf uur wiskunde wordt het ontbinden nog verder uitgebreid in het vierde jaar bij veeltermen. Als voorbereiding op de meer gevorderde wiskunde van de derde graad kan in die richtingen in het vierde jaar dan expliciet geoefend worden op het rekenen met rationale lettervormen. Hierbij speelt ontbinden in factoren opnieuw een grote rol. 

We pleiten ervoor om daarnaast de geoefende vaardigheden zo vaak mogelijk zijdelings aan bod te laten komen in de leerstof van de tweede graad doorheen dit proces. Denk bijvoorbeeld aan tweedegraadsvergelijkingen met parameters of het bewijzen van goniometrische identiteiten met oefeningen waarin een rationale vorm of een ontbinding met een merkwaardig product voorkomt. 

In de derde graad worden deze vaardigheden ingezet bij rationale functies, goniometrie, complexe getallen, afgeleiden en verloop van functies, integralen\ldots Het is dan belangrijk om de algebraïsche vaardigheden vast te zetten die in de tweede graad werden aangeleerd en ingeoefend. Aan het einde van de derde graad zijn leerlingen op die manier optimaal voorbereid op het hoger onderwijs en kunnen ze vlot de geschikte rekenprocedures selecteren en toepassen bij het oplossen van problemen.


\subsection{Uitbouwen technieken in het vierde jaar}

Omdat dit niet meer expliciet aan bod komt in de handboekenreeksen werken we hier een aantal lesactiviteiten uit om de verschillende rekentechnieken in te oefenen. We geven bij elke techniek eerst een uitgewerkt voorbeeld, waarna we de verschillende stappen in de procedure expliciteren. Daarna geven we een selectie van een aantal voorbeeldoefeningen.

Bij de uitgewerkte voorbeelden kun je een aantal discussievragen stellen aan de leerlingen om zo samen tot de explicitering van de procedure te komen. Leerlingen oefenen daarbij ook het gebruik van correcte vakterminologie. Voorbeelden van zo'n discussievragen:
\begin{itemize}
	\item Welke stappen komen aan bod bij het oplossen van het probleem? Waarom werken ze in deze volgorde? Zouden ze ook in een andere volgorde kunnen werken?
	\item Kon het probleem worden opgelost met minder stappen?
	\item Kun je een andere manier bedenken om dit probleem op te lossen?
	\item Zal deze strategie altijd werken? Waarom?
	\item Bedenk een ander probleem waarbij deze strategie werkt.
	\item Kan je iets aan de opgave veranderen zodat je dezelfde strategie niet meer kunt toepassen?
	\item Hoe kun je de uitwerking nog aanpassen zodat ze duidelijker wordt voor anderen?
	\item Welke al geziene technieken komen voor in de uitwerking? (Bijvoorbeeld: ontbinden in factoren.)
\end{itemize}


De opsomming van de verschillende stappen in het stappenplan is minder belangrijk en moet zeker niet uit het hoofd geleerd worden door leerlingen. 


\end{multicols}
 \afbeelding[0.85\textwidth]{img/hordelopen.jpg}{}{}

 \vspace{0.5cm}
\begin{lesactiviteit}{Rekenen met rationale lettervormen}
	\begin{kader}{}{Vereenvoudigen van rationale lettervormen: uitgewerkte voorbeelden}
		$\displaystyle \frac{a^2 b}{5ab^3}=\frac{a}{5b^2}$\\
		
		$\displaystyle \frac{a-b}{a^2-b^2}=\frac{a-b}{(a-b)(a+b)}=\frac{1}{a+b}$\\
		
		$\displaystyle \frac{-x^3+x^2+2x-2}{x^2-1}=\frac{(x-1)(-x^2+2)}{(x-1)(x+1)}=\frac{-x^2+2}{x+1}$\\
\end{kader}

\textbf{Stappenplan voor het vereenvoudigen van rationale lettervormen}
\begin{enumerate}
	\item Ontbind eerst de teller en de noemer in factoren.
	\item Deel vervolgens teller en noemer door hun gemeenschappelijke factoren.
\end{enumerate}
\vspace{1cm}

\begin{kader}{}{Optellen van rationale lettervormen: uitgewerkte voorbeelden}
Voorbeeld 1:
\begin{align*}
\frac{-1}{x^2+2x}+\frac{2}{x^2-4}&=\frac{-1}{x(x+2)}+\frac{2}{(x-2)(x+2)} \\
&=\frac{-1 (x-2)}{x(x+2)(x-2)}+\frac{2x}{(x-2)(x+2)x}\\
&=\frac{-x+2+2x}{x(x-2)(x+2)} \\
&=\frac{(x+2)}{x(x-2)(x+2)} \\
&=\frac{1}{x(x-2)}
\end{align*}

Voorbeeld 2:
\begin{align*}
	\frac{1}{a-b}+\frac{3}{b-a} & = \frac{1}{a-b}+\frac{-3}{a-b}\\
								& = \frac{1-3}{a-b}\\
								& = \frac{-2}{a-b}
\end{align*}
\end{kader}	
\textbf{Stappenplan voor het optellen van rationale lettervormen}
\begin{enumerate}
	\item Ontbind de noemers in factoren.
	\item Zet de breuken op gelijke noemer: dit is het kleinste gemeen veelvoud van de noemers.
	\item Tel de tellers op en behoud de noemers. (De noemers blijven standaard ontbonden.)
	\item Vereenvoudig de breuk, als dit mogelijk is.
\end{enumerate}
\vspace{1cm}


\begin{kader}{}{Vermenigvuldigen van rationale lettervormen: uitgewerkt voorbeeld}
	\begin{align*}
		\frac{a^2+2ab+b^2}{a^2-ab}\cdot \frac{3a^2}{2a+2b}&=\frac{(a^2+2ab+b^2)\cdot(3a^2)}{(a^2-ab)\cdot (2a+2b)}\\
		&=\frac{(a+b)^2\cdot 3a^2}{a(a-b)\cdot 2(a+b)}\\
		&=\frac{3a(a+b)}{2(a-b)}
	\end{align*}
\end{kader}	
\textbf{Stappenplan voor het vermenigvuldigen van rationale lettervormen}
\begin{enumerate}
	\item Vermenigvuldig de tellers met elkaar en de noemers met elkaar.
	\item Vereenvoudig de breuk, indien mogelijk.
\end{enumerate}
\vspace{1cm}
		

\begin{kader}{}{Delen van rationale lettervormen: uitgewerkt voorbeeld}
	\begin{align*}
		\frac{(a+3)}{1+\frac{a}{a-1}}&=\frac{a+3}{\frac{a-1+a}{a-1}}\\
		&=\frac{a+3}{\frac{2a-1}{a-1}}\\
		&=(a+3)\cdot \frac{a-1}{2a-1}\\
		&=\frac{(a+3)(a-1)}{2a-1}		
	\end{align*}
\end{kader}	
\textbf{Stappenplan voor het delen van rationale lettervormen}
\begin{enumerate}
	\item Schrijf de opgave als een quotiënt van twee breuken.
	\item Maak gebruik van de rekenregel: $\displaystyle \frac{\frac{a}{b}}{\frac{c}{d}}= \frac{a}{b} \cdot \frac{d}{c}$.
	
	Let op: je mag dit enkel doen als de teller en de noemer allebei een breuk zijn, en bijvoorbeeld niet als de noemer een som van twee breuken is.
\end{enumerate}
\end{lesactiviteit}

\begin{lesactiviteit}{Een selectie van oefeningen (deel 1)}
\begin{multicols}{2}
\begin{vraagenantwoord}
    \vraag{Vereenvoudig indien mogelijk.}
    \begin{itemize}[label=\textcolor{\rubriekkleur}{\Large\textbullet}] 
    \item $\frac{x^2-y^2}{x-y}$
    \item $\frac{x^2+2x+1}{x^2-1}$
    \item $\frac{2x^2-5x+3}{-4x^2+8x-3}$
    \end{itemize}
    \vraag{Pas eerst (kruisgewijs) vereenvoudigen toe vooraleer de breuken te vermenigvuldigen.}
    \begin{itemize}[label=\textcolor{\rubriekkleur}{\Large\textbullet}] 
    \item $\frac{3(a-b)}{a(a+b)} \cdot \frac{a^2}{9(a-b)}$
    \item $\frac{(x+1)(x-2)}{x^2(x+3)} \cdot \frac{x^2}{(x-3)(x+1)}$
    \end{itemize}
    \vraag{Vereenvoudig indien mogelijk.}
    \begin{itemize}[label=\textcolor{\rubriekkleur}{\Large\textbullet}] 
    \item $\frac{(b-a)^3}{a^2-b^2}$
    \item $\frac{x^2+1}{x^4+2x^2+1}$
    \item $\frac{(2x-y)^2-(x+3y)^2}{(3x-y)^2-(2x+3y)^2}$
    \end{itemize}
    \vraag{Vereenvoudig indien mogelijk. Zonder eerst in de teller en noemer de gemeenschappelijke factoren af.}
    \begin{itemize}[label=\textcolor{\rubriekkleur}{\Large\textbullet}] 
    \item $\frac{3x^4-3x^2}{x^4+x^2}$
    \item $\frac{a^2-ab}{b^2-ab}$
    \item $\frac{2(x-2)+(x-2)x^2}{x-2}$
    \end{itemize}
    %\vraag{Bereken en vereenvoudig indien mogelijk.}
   %\begin{itemize}[label=\textcolor{\rubriekkleur}{\Large\textbullet}] 
   % \item $\frac{3}{x+1} + \frac{2x+1}{x-2}$ 
    %\item $\frac{4a+1}{a+2} + \frac{5a}{3(a+2)}$
    %\item $1-\frac{x-2}{3x+7}$
    %\end{itemize}
\end{vraagenantwoord}
 \end{multicols}   
\end{lesactiviteit}








%\afbeelding[0.9\linewidth]{img/LAlgebra/LEls/algebra_oef1.jpg}{Een selectie van oefeningen, deel 1}{fig:oef1}

\begin{multicols}{2}

De lesactiviteiten 'Een selectie van oefeningen deel 1 en deel 2' en bieden een selectie van oefeningen uit een bundel die de auteurs van Wiskundeplan ter beschikking stellen op \url{www.wiskundeplan.be}. Wanneer leerlingen fouten maken tegen een bepaalde techniek of wanneer ze twijfelen of ze een bepaalde regel mogen toepassen, kun je hen als tip geven dat ze het moeten proberen met concrete getallen. Numerieke voorbeelden zijn handig om fouten te ontkrachten.

Het is belangrijk om voldoende lang basisoefeningen te maken (meer dan hieronder zijn opgesomd) en niet te snel over te stappen op meer ingewikkelde lettervormen. Inspiratie voor oefeningen kun je vinden in oude handboeken waarin dit topic nog aan bod kwam, in de bundel van Wiskundeplan, in zomercursussen van universiteiten zoals bijvoorbeeld in het onderdeel elementair rekenen van het zelfstudiepakket wiskunde (informatica) van de KU Leuven, \ldots

\end{multicols}

\begin{lesactiviteit}{Een selectie van oefeningen (deel 2)}
\begin{multicols}{2}
\begin{vraagenantwoord}[start=5]
   \vraag{Bereken en vereenvoudig indien mogelijk.}
   \begin{itemize}[label=\textcolor{\rubriekkleur}{\Large\textbullet}]
    \item $\frac{3}{x+1} + \frac{2x+1}{x-2}$ 
    \item $\frac{4a+1}{a+2} + \frac{5a}{3(a+2)}$
    \item $1-\frac{x-2}{3x+7}$
    \end{itemize}
    \vraag{Vereenvoudig indien mogelijk.}
    \begin{itemize}[label=\textcolor{\rubriekkleur}{\Large\textbullet}] 
    \item $\frac{2a^2}{(a-1)(a+1)}- \frac{a^2+1}{a^2-1}$
    \item $\frac{1}{(a-b)^2+ \frac{1}{(b-a)^2}}$
    \item $\frac{1}{x-1} + \frac{x+1}{(x-1)(x-3)}$
    \item $\frac{8}{(2+x)(x-2)} + \frac{2}{2-x}$
    \end{itemize}
    \vraag{Bereken en vereenvoudig indien mogelijk.}
    \begin{itemize}[label=\textcolor{\rubriekkleur}{\Large\textbullet}] 
    \item $\frac{x+6}{x^2-16} - \frac{x+1}{x^2-4x}$
    \item $\frac{x+1}{x-1} + \frac{x-1}{x+1} - \frac{4}{x^2-1}$
    \item $\frac{4x}{x^2-4} + \frac{2}{x^2-5x+6}$
    \end{itemize}
    \vraag{Bereken en vereenvoudig indien mogelijk.}
    \begin{itemize}[label=\textcolor{\rubriekkleur}{\Large\textbullet}] 
    \item $\frac{x^2-x}{x-1} \cdot \frac{5x-5}{1-x}$
    \item $\frac{x^3-8}{x^3+8} \cdot \frac{x^2+4x+4}{x^2-4}$
    \item $\frac{xy+y^2}{(x-y)^2} \cdot \frac{x^2-xy}{(x+y)^2}$
    \end{itemize}
    \vraag{Bereken en vereenvoudig indien mogelijk.}
    \begin{itemize}[label=\textcolor{\rubriekkleur}{\Large\textbullet}] 
    \item $(1-\frac{1}{x-2})(x-3+\frac{1}{x-3})$
    \item $(x-\frac{x+2}{x})(5-\frac{1}{x+1})$
    \end{itemize}
     \vraag{Bereken en vereenvoudig indien mogelijk.}
    \begin{itemize}[label=\textcolor{\rubriekkleur}{\Large\textbullet}] 
    \item $\frac{\frac{5x^2-5}{x^2}}{\frac{6x+6}{x^3}}$
    \item $\frac{3+\frac{1}{x-2}}{3-\frac{1}{x-5}}$
    \end{itemize}
     \vraag{Bereken en vereenvoudig indien mogelijk.}
    \begin{itemize}[label=\textcolor{\rubriekkleur}{\Large\textbullet}] 
    \item $(\frac{1}{a}+\frac{1}{b})^{-1}$
    \item $(\frac{a}{a-b}-\frac{b}{b-a})^{-1}$
    \end{itemize}
     \vraag{Bereken en vereenvoudig indien mogelijk.}
    \begin{itemize}[label=\textcolor{\rubriekkleur}{\Large\textbullet}] 
    \item $\frac{(a-1)^3-(a+1)(a-1)^2}{(a-1)^6}$
    \item $\frac{3x+3}{x^2-2x}-\frac{2x+2}{x^2-4}$
    \item $(\frac{1}{y^2}-\frac{1}{x^2})\cdot \frac{y}{1-\frac{y}{x}}$
    \item $\frac{x^2-y^2+x+y}{x+y}$
    \item $\frac{1+x}{1+2x+x^2}-\frac{1+x}{1+x^3}$
    \item $\frac{1}{\frac{a-b}{ab}}\cdot (\frac{b}{a}-\frac{a}{b})$
    \end{itemize}
\end{vraagenantwoord}
 \end{multicols}   
\end{lesactiviteit}



%\afbeelding[0.9\linewidth]{img/LAlgebra/LEls/algebra_oef2.jpg}{Een selectie van oefeningen, deel 2}{fig:oef2}
%\afbeelding[0.9\linewidth]{img/LAlgebra/LEls/algebra_oef3.jpg}{Een selectie van oefeningen, deel 3}{fig:oef3}

\begin{multicols}{2}
	\section{ICT voor remediëring}
    Tot nog toe belichtten we belangrijke algebraïsche rekentechnieken voor de tweede graad, namelijk ontbinden in factoren en werken met letters in breuken.
    
    We nemen nu een kijkje naar hoe je ICT kunt inzetten om je algebralessen te ondersteunen, zowel voor de tweede als derde graad.
    De nadruk zal vooral liggen op de technologische mogelijkheden, al geven we soms een didactische suggestie.
    Verschillende `tools' worden besproken als voorbeeld bij fictieve vragen van leerlingen zoals ``Kunt u de tussenstapjes tonen?'' of ``Heeft u nog andere oefeningen?''.     
    We pogen hier zeker geen exhaustieve lijst aan te bieden, maar hopen je een startpunt te bieden om zelf op zoek te gaan.
    
    Je zult verder merken dat er regelmatig hyperlinks in de tekst staan, wat eigen is aan de aard van dit onderwerp.
    We hebben ervoor gezorgd dat die hyperlinks telkens bij een kernwoord horen zodat lezers op papier gemakkelijk zelf online kunnen zoeken.
    Op het einde van deze loep vind je ook een lijst met alle hyperlinks uit dit deel.

    Voor de vraaggerichte paragraafjes, beginnen we eerst met een algemene reflectie op welke eigenschappen gewenst zijn bij een ICT-tool.

    \subsection{Gewenste functionaliteiten} 


    
    We lijsten in tabel~\ref{tab:functionaliteiten} enkele functionaliteiten op, die je kunnen helpen bij je eigen zoektocht naar interessante ICT-oplossingen om de leerlingen te ondersteunen. 

        In wat volgt zullen we niet elk van de bovenstaande kenmerken bespreken bij iedere tool, om je als lezer niet te overbelasten.
    Houd het echter in je achterhoofd, zowel in je eigen zoektocht als tijdens het verder lezen van deze loep.

    \end{multicols}

        \begin{tabel}
		\begin{tabular}{m{2cm}|m{13cm}}  % tekstbreedte=7.5cm, elke tabelkolom = breedte + 2x1mm marge (2x want marge links & rechts)
			\hlinethick  % dikkere lijn, gebruik deze rond de hoofding (eerste rij) en om de tabel te eindigen
            \tableheader{Eigenschap} & \tableheader{Uitleg}\\
            Prijs & Bij voorkeur, gratis. Eventueel deel van een reeds bestaand schoolabonnement.\\
            Identificatie & Als de leerlingen geïdentificeerd worden in de tool, kun je ze op de lange termijn volgen/individuele tips geven/\ldots Als ze anoniem oefenen, verlaagt dat de oefendrempel voor hen.\\
            Adaptiviteit & Zijn de verschillende oefeningen gerangschikt volgens moeilijkheidsniveau? En krijgt een leerling automatisch oefeningen van het juiste moeilijkheidsniveau?\\
            Stap-voor-stap & Krijgen de leerlingen de tussenstappen te zien? En is er aandacht voor alternatieve oplossingsmethodes?\\
            Fine-tunen & Kun je specifieke algebraregels oefenen, of enkel grotere gemengde oefeningen?\\
            Correctheid & Maakt de tool fouten?\\
            \hlinethick
		\end{tabular}
	\caption{Functionaliteiten van software.}\label{tab:functionaliteiten}
	\end{tabel}

\begin{multicols} {2}
    

    
    \subsection{Kunt u de tussenstapjes tonen?}
    De meest elementaire vorm van een ICT-hulpmiddel voor algebra, is software om het eindantwoord te controleren, en bij voorkeur om ook de tussenstapjes te zien.

   
    
    In \autoref{fig:voorbeeldtussenstappen} werkt de computer het voorbeeld \(\left(x+4\right)\left(x+6\right)\) met tussenstappen uit.
    We gaven zelf ``\(x+1\)'' puur als voorbeeld van een fout antwoord.

 ``Symbolisch rekenen'' door computers bestaat al decennia lang, en je kunt je dus voorstellen dat er al een overvloed aan ICT-tools bestaat die kan wat \autoref{fig:voorbeeldtussenstappen} toont.
    Dit is effectief het geval, maar we pikken er specifiek ``\href{https://www.wolframalpha.com/input?i=factor+x%5E2-y%5E2}{WolframAlpha}'' uit.
    Dit is een gebruiksvriendelijke en gratis app/website

    
    \afbeelding{img/LAlgebra/LAlgebraICT/voorbeeldje-tussenstapjes.png}{Digitale tussenstappen voor algebra}{fig:voorbeeldtussenstappen}

     

    die ook de tussenstapjes laat zien.
    Indien je berekening te veel tijd vraagt, zal de website je echter niet alles tonen en om betaling vragen.
    In dat geval kun je gebruik maken van een gratis, afgeleide versie genaamd ``WolfreeAlpha''.
    Als je deze term opzoekt in een zoekmachine, zul je \href{https://wolfreealpha.onrender.com/}{verschillende websites} tegenkomen die zonder betaling alle tussenstappen tonen.
    De beschikbare websites van WolfreeAlpha veranderen regelmatig; een up-to-date overzicht vind je \href{https://www.reddit.com/r/GetStudying/comments/x4luun/get_wolframalpha_pro_feature_for_free/}{online op discussieforum ``Reddit''}.
    
    Er bestaan ook apps die je op je smartphone kunt installeren.
    Een goed voorbeeld is \href{https://photomath.com/}{Photomath}.
    Met die app kun je simpelweg de opgave inscannen met de camera, en vervolgens genereert de app de oplossing met tussenstapjes.
    
    Dit zijn slechts twee tools als voorbeeld; je kunt er nog veel meer terugvinden door online te zoeken naar ``Computer Algebra System with steps''.
    Maar \textbf{hoe zet je die tools dan in?}
    Dergelijke basis-tools kun je het gemakkelijkst inzetten als verbetersleutel voor de leerlingen.
    Zij kunnen op die manier zelfstandig (minimale) feedback op hun antwoordpoging krijgen.
    We raden je wel aan om duidelijk te communiceren naar je leerlingen dat deze tools soms een andere oplossingsmethode gebruiken.
    Het antwoord van de leerling hoeft dus niet fout te zijn als het verschilt van de getoonde tussenstappen.
    Meer zelfs, voor dergelijke verschillen kan de leerkracht gedetailleerde feedback geven, en misschien interessante discussies openen die tot meer inzicht leiden.

    \subsection{Kunt u die techniek nog eens uitleggen?}
    Naast oefenmogelijkheden zijn er ook digitale cursussen te vinden, die de algebraïsche technieken van de grond opbouwen.
    Je kunt dergelijk materiaal gebruiken om te differentiëren tussen leerlingen: extra herhaling om de fundamentele technieken te oefenen of uitdaging voor geavanceerdere algebraïsche regels.
    Voordeel is dat de leerling op zijn eigen tempo kan werken, mits een kwaliteitscontrole van de cursus door de leerkracht.

    Het leerplatform van het gebruikte handboek is een eerste optie.
    Dat biedt vaak al veel mogelijkheden en is ook mooi geïntegreerd in het lesmateriaal dat leerlingen gewoon zijn.
    Het nadeel hiervan is vaak dat de tussenstappen niet, onvoldoende of eenzijdig zijn uitgewerkt.

    Er bestaan ook veel video's op YouTube die wiskundige concepten introduceren of herhalen.
    In het Nederlands heb je bijvoorbeeld de \href{https://www.youtube.com/watch?v=Utx1apbbuao}{WiskundeAcademie} of \href{https://www.youtube.com/watch?v=pjSZ2tvGbvA}{Math with Menno}.
    Ook \href{https://youtu.be/rR4FfB-FZJw?si=K42IUlAAXmSWXuyL}{KhanAcademy} is een kwaliteitsvol YouTube-kanaal (in het Engels). 
    Meer zelfs, deze non-profit organisatie werd opgericht in 2009 en is met de jaren uitgegroeid tot een gigantisch project dat gratis, online leren promoot met meer dan 8000 video's op YouTube en een volledig bijbehorend leerplatform.
    Op \href{https://www.khanacademy.org/math/algebra-basics/alg-basics-quadratics-and-polynomials}{het leerplatform} kun je een gratis account maken, en door gestructureerde modules lopen met toegevoegde oefeningen.
    Deze website is grotendeels in het Engels, maar hij wordt stelselmatig omgezet naar het Nederlands.
    
    Tot slot vermelden we de voorbereidingscursussen van de Vlaamse universiteiten voor beginnende bachelorstudenten.
    Bijvoorbeeld KU Leuven heeft \href{https://set.kuleuven.be/voorkennis/zomercursus/zomercursusZ/rekenen/breuken}{een volledig online cursus} met toegevoegde oefeningen, die ook de algebraïsche rekentechnieken behandelen.
    Ook UHasselt heeft een gratis beschikbare cursus, in \href{https://sites.google.com/uhasselt.be/platformwiskunde/leerplatform/leerboek/algebra-en-getaltheorie/veeltermen/veeltermen-tot-graad-2}{pdf-vorm online}.
    UGent vermeldt de cursus ``\href{https://elemath.ugent.be/}{EleMath}'', maar we vinden enkel online oefeningen in een eigen leerplatform (waar je ook anoniem kunt aanmelden).
    De eigenlijke cursus lijkt enkel apart verkrijgbaar.
    Tot slot heeft UAntwerpen ook zijn eigen leerplatform ``Aan De Slag'' waar iedereen gratis een account kan maken.
    In dat platform vind je \href{https://www.uantwerpen.be/nl/projecten/aan-de-slag/zelfstudiepakketten/}{het zelfstudiepakket} ``Wiskunde: Voorkennis en opfrissing voor alle opleidingen''.
    Dit bestaat uit een theoriecursus in pdf, video's die de technieken voortonen en meerkeuze-oefeningen om zelf te oefenen.
    Er lijken geen digitale oefeningen bij te horen waar de leerling feitelijk een algebraïsche uitdrukking moet ingeven in een tekstvak.

    \subsection{Heeft u nog andere oefeningen?}
    De volgende tools genereren automatisch nieuwe oefeningen.
    Merk op: dit gaat verder dan een manueel bijgehouden lijst van oefeningen: het onderliggende ``Computer Algebra System'' (kortweg CAS) genereert zelf nooit-eerder-geziene oefeningen.
    Als de leerlingen die oplossen krijgen ze eventueel feedback tussendoor, en automatische feedback bij het indienen van de oplossing.
    Als ze meer oefening nodig hebben, kunnen ze simpelweg meer oefeningen vragen met de klik van een knop.

    Als eerste presenteren we \href{https://practice.algebrakit.com/nl-BE/overview}{Algebrakit}.    
    Op deze website kun je gratis oefenthema's selecteren, die samenstellen in een oefenbundel en vervolgens de link aan je leerlingen bezorgen.
    De specifieke link houdt bij wat de leerling heeft geprobeerd en dat kun je inkijken zonder accounts voor je leerlingen aan te maken.
    Wanneer zij op de link komen, krijgen ze automatisch gegenereerde oefeningen te zien met de mogelijkheid voor een hint.
    Ze kunnen hun eindantwoord en tussenstappen controleren en om nieuwe oefeningen vragen \textbf{binnen verschillende moeilijkheidsniveaus}.
    In \autoref{fig:demoalgebrakit} zie je wat een leerling te zien krijgt.
    \afbeelding{img/LAlgebra/LAlgebraICT/demo-algebrakit.png}{Voorbeeld Algebrakit}{fig:demoalgebrakit}
    Je zult merken dat er weinig onderwerpen beschikbaar zijn in het Nederlands, maar dat kun je verhelpen door ``United States'' als regio te kiezen.
    In deze oefeningen heeft de taal weinig invloed, dus het is nog steeds een interessante tool om te bekijken.
    
    We merken op dat het factoriseeralgoritme van Algebrakit niet altijd de meest gefactoriseerde vorm geeft.
    Volgens de website is \(3s^3t+3s^2t-27st\) namelijk volledig gefactoriseerd in 
    \[
    3st\left(s^2+s-9\right)
    \]
    Als je de discriminant controleert van de tweedegraadsfactor, zie je echter dat er nog factoren uit te splitsen zijn, namelijk:
    \[
    3st\left(s-\frac{-1+\sqrt{37}}{2}\right)\left(s-\frac{-1-\sqrt{37}}{2}\right)
    \]
    Hoewel het een heel rijk platform is, lijkt het onderliggende CAS dus niet volledig.
    Dit gebrek kan echter een ideale leeropportuniteit zijn door dit te analyseren met je leerlingen.    
    Wij vermoeden alleszins dat Algebrakit slechts één ontbindingsregel per oefening kan toepassen en daarom niet de volledige factorisatie vindt.
    
    Op \href{https://quickmath.com/webMathematica3/quickmath/algebra/simplify/basic.jsp}{QuickMath.com} kun je ook willekeurige oefeningen laten genereren en de tussenstappen zien.
    De twee algebraïsche vaardigheden die we in deze loep behandelen zijn gemakkelijk terug te vinden: ontbinden in factoren bij ``Factor'', werken met letters in breuken is deel van ``Simplify''.
    Opgepast, de gratis versie van deze tool toont niet alle tussenstappen in de modeloplossing.
    Je kunt in dat geval gewoon de opgave fotograferen met PhotoMath of teruggrijpen naar WolframAlpha.
    Merk tot slot ook op dat QuickMath geen moeilijkheidsinstelling heeft, en in onze experimenten vrij pittige oefeningen genereert.
    Wij kwamen bijvoorbeeld de breuk
    \[
    \frac{\frac{a^3 - b^3}{x + y}}{\frac{a^2 + a b + b^2}{x + y}}
    \]
    tegen bij ``Simplify''.
    In een eerdere oefening stelden we reeds een oefening voor waar leerlingen de efficiëntie beoordelen van verschillende algebraïsche afleidingen (zie \autoref{q:efficientie}).
    Deze pittige oefeningen van QuickMath zijn ideaal om hetzelfde te doen, maar met computeralgoritmes.
    In de bovenstaande uitdrukking vereenvoudigen wij mensen namelijk onmiddellijk tot \(\frac{a^3-b^3}{a^2 + a b + b^2}\) door de factor \(x+y\) te schrappen.
    WolframAlpha stelt echter iets inefficiënter voor als eerste stap, namelijk om met het omgekeerde van de noemer te vermenigvuldigen:
    \[
    \ldots = \frac{a^3 - b^3}{x + y}\cdot\frac{x+y}{a^2 + a b + b^2} = \ldots
    \]
    Door dergelijke verschillen te bestuderen oefen je het inzicht van je leerlingen, en maak je nog maar eens duidelijk dat de computer slechts één oplossingsmethode toont.

    We vermeldden reeds Algebrakit en QuickMath.
    Ook in WolframAlpha kun je oefeningen genereren in hun ``\href{https://www.wolframalpha.com/problem-generator/?scrollTo=Algebra}{Problem Generator}'' waar je uit een hele lijst topics kan kiezen.
    In tegenstelling tot QuickMath, is ontbinden in factoren en letterrekenen in breuken hier geen apart onderwerp.
    Je zult dus zelf wat moeten puzzelen.
    We vermelden kort dat WolframAlpha ook \href{https://www.wolframalpha.com/pro/pricing/educators/}{een betalende versie voor leerkrachten} heeft, waarmee je oefeningenblaadjes kunt genereren en afdrukken om uit te delen.
    Met die betalende versie kun je ook altijd alle tussenstapjes zien, zonder gebruik te moeten maken van iets als WolfreeAlpha.
    
    Een platform om mee te eindigen in dit deel, is het Britse platform ``\href{https://uk.ixl.com/maths/algebra}{IXL}''.
    Hier kan men willekeurige oefeningen genereren, een voorbeeld als hint vragen en het juiste antwoord met tussenstappen zien.
    Je kunt door een groot deel van het Britse wiskundecurriculum neuzen, en alle oefeninggeneratoren zijn volgens onderwerp terug te vinden.
    De algebraïsche technieken van deze loep vind je bij ``Year 11''.
    Ontbinden in factoren staat bij ``\href{https://uk.ixl.com/maths/year-11/factorise-by-grouping}{factorize by grouping}'' en werken met letters in breuken bij ``\href{https://uk.ixl.com/maths/year-11/simplify-rational-expressions}{simplify rational expressions}''.
    De kracht van deze website komt naar boven bij de betalende versie: dan kunnen leerlingen accounts maken, en kun je in detail hun vooruitgang opvolgen en aparte oefeningen voorstellen.

    Alles samengenomen, zijn er heel veel ICT-tools terug te vinden voor algebraonderwijs.
    Op \href{https://www.homeschoolmath.net/online/curricula.php}{homeschoolmath.net} vind je een lijst van andere tools die hier niet aan bod kwamen.
    Wellicht vind je er nog meer als je zelf online rondzoekt.

    We eindigen deze uitstap naar ICT voor algebra-onderwijs met enkele uitdagende topics: zelf automatische algebravragen programmeren en het inzetten van AI\ldots

    \subsection{Heeft u die vraag zelf gemaakt?!}
    We raakten al eerder Computer Algebra Systemen aan, en zoals gezegd bestaan er heel veel programmeertalen die symbolisch rekenen ondersteunen.
    Je kunt je voorstellen dat alles wat we tot hiertoe bekeken hebben, in Python geprogrammeerd kan worden.
    Zoiets zelf maken, zou een huzarenwerkje worden.
    Het gemakkelijkste is om een bestaand vragenplatform te gebruiken dat al met algebraïsche antwoorden overweg kan.
    Dat betekent dat leerlingen een antwoord als \(x+1\) kunnen ingeven, en het platform begrijpt dat dit hetzelfde is als \(2x+1-x\).
    Zo bestaan er verschillende, maar we focussen hier op \href{https://www.numbas.org.uk/}{NUMBAS}.
    
    NUMBAS is een gratis, open-source vragenplatform van de Britse universiteit van Newcastle, en het focust specifiek op wiskunde en andere wetenschappen.
    Meer nog, je kunt vrij gemakkelijk zelf een vraag maken die telkens een andere veelterm genereert en die de leerlingen dan moeten factoriseren/vereenvoudigen.
    Neem maar eens een kijkje op hun \href{https://numbas.mathcentre.ac.uk/exam/1973/numbas-website-demo/embed/}{demo-website}.
    Om zelf van start te gaan kun je de uitleg volgen op \href{https://docs.numbas.org.uk/en/latest/tutorials/write-question.html}{hun documentatiewebsite}.
    Toegegeven, er is een redelijke leercurve aan verbonden.
    Er is echter veel voorbeeldmateriaal beschikbaar, ontwikkeld door \href{https://numbas.mathcentre.ac.uk/explore/}{verschillende Britse universiteiten} en door andere leerkrachten van over heel de wereld.
    Je kunt gemakkelijk rondneuzen volgens onderwerp, en vervolgens een vraag kopiëren die jij ietwat wilt aanpassen of vertalen.
    Op die manier hoef je zelf geen programmeerexpert te worden, maar kun je stap-voor-stap iets bijleren over een specifieke functionaliteit van NUMBAS en kunnen je leerlingen op specifieke algebraïsche technieken oefenen.

    We werken even een voorbeeld uit.
    Bij \href{https://numbas.mathcentre.ac.uk/project/601/browse/Foundation%2520mathematics/}{de vragen van de universiteit van Newcastle}, vinden we \href{https://numbas.mathcentre.ac.uk/question/12000/combining-algebraic-fractions-1/preview/}{een vraag} waar twee letterbreuken opgeteld moeten worden (zie \autoref{fig:numbasvoorbeeld}).
    
    \afbeelding{img/LAlgebra/LAlgebraICT/voorbeeldnumbas.png}{Optellen van twee letterbreuken in NUMBAS}{fig:numbasvoorbeeld}
    Stel dat je graag de willekeurige breuken wat had aangepast, zodanig dat altijd het merkwaardig product \(A^2-B^2\) kan gebruikt worden voor de gemeenschappelijke noemer.
    Op die manier combineer je voor de leerlingen zowel ontbinden in factoren als werken met letterbreuken.
    We maakten daarvoor een kopie van deze vraag naar ons eigen NUMBAS-account en vertaalden al de tekst naar het Nederlands.
    Bij het opstarten van de originele vraag werden willekeurige coëfficiënten \(a\), \(b\), \(c\) en \(d\) gekozen in
    \[
    \frac{a}{x+b} + \frac{c}{x+d}
    \]
    In onze versie (die je \href{https://numbas.mathcentre.ac.uk/question/155648/voorbeeld-uitwiskeling/}{hier online} terugvindt) moeten we die opgave licht aanpassen om het merkwaardig product te kunnen gebruiken.
    Daarom veranderden we bij het ``Parts''-tabblad naar de onderstaande opgave in \autoref{fig:aanpassingopgaveNUMBAS}.
    \afbeelding[0.35\textwidth]{img/LAlgebra/LAlgebraICT/opgaveaanpassingNumbas.png}{Aangepaste opgave in NUMBAS}{fig:aanpassingopgaveNUMBAS}
    Zoals je kunt zien, zal deze vraag altijd oplosbaar zijn met het merkwaardig product \(A^2-B^2\) voor de gemeenschappelijke noemer.
    Je kunt de vraag vervolgens doorsturen naar je leerlingen om te oefenen via een link.
    Ook kun je het delen met je collega's en er samen verder aan werken.

    Indien je geïnteresseerd bent, neem eens een kijkje naar ons voorbeeld en het documentatiemateriaal.
    Het is aan jou om te kiezen hoe ver je hierin gaat.
    Onze voorbeeldvraag controleert bijvoorbeeld nog niet of de leerling effectief maar één breuk ingeeft, wat een interessante toevoeging zou zijn om te programmeren.
    We merken tot slot op dat NUMBAS slechts één platform is met algebraïsche antwoorden.
    Een andere populaire is \href{https://stack-assessment.org/}{STACK}, ontwikkeld door ETH Zürich (Zwitserland), met een iets rijker onderliggende CAS maar minder beschikbaar materiaal.

    \subsection{Gaat dat ook met AI?}
    Laat het ons simpelweg eens uitproberen als inleiding op dit onderwerp.
    We stelden aan \href{https://chatgpt.com/}{ChatGPT 3.5} de volgende vraag:
    \[
    \text{How do I factor }\left(a+b-c\right)^2 - \left(a-b+c\right)^2\text{ ?}
    \]
    en kregen het onderstaande antwoord (vertaald in het Nederlands).
    \begin{kader}{ChatGPT}{Antwoord}
        \begin{minipage}{\textwidth}
        Om \(\left(a+b-c\right)^2 - \left(a-b+c\right)^2\) te ontbinden in factoren gebruiken we het merkwaardig product \(a^2-b^2=(a-b)(a+b)\).\newline
        Toegepast op de gegeven uitdrukking levert dat
        \[
        \begin{array}{l}
        \left(a+b-c\right)^2 - \left(a-b+c\right)^2 =\\ 
        \left((a+b-c)+(a-b+c)\right) \cdot\\
        \left((a+b-c)-(a-b+c)\right)= \\
            2a\left(2b-2c\right)=\\
            4a(b-c)
        \end{array}
        \]
        \end{minipage}
    \end{kader}
    Laat ons echter onmiddellijk eens kijken naar een lichtjes andere vraag:
    \[
    \text{Please factor }\left(a+b-c\right)^2 - \left(a-b+c\right)^2
    \]
    Dit is wiskundig gezien exact dezelfde vraag.
    We kregen echter een verschillend antwoord als we deze vraag stelden aan ChatGPT 3.5 (in een ander gesprek).
    \begin{kader}{ChatGPT}{Antwoord}
        Om deze uitdrukking te ontbinden kunnen we het merkwaardig product \(a^2 - b^2 = (a + b)(a - b)\) gebruiken.\newline
        Als we dit toepassen op onze uitdrukking, krijgen we:
        \[
        \begin{array}{l}
        \left(a+b-c\right)^2 - \left(a-b+c\right)^2 =\\ 
        \left((a+b-c) + (a-b+c)\right)\cdot\\
        \left((a+\textcolor{red}{b}-c) - (a-\textcolor{red}{b}+c)\right)=\\
        \left(2a\right)\left(2c\right)=\\
        4ac
        \end{array}
        \]
    \end{kader}
    Je merkt het hopelijk op: ChatGPT heeft een fout gemaakt en is een \(b\) verloren in de berekening.
    En dat puur door een lichtjes andere vraag te stellen!
    Het is ook vreemd bij beide antwoorden om de variabele \(a\) zowel voor een variabele op zich, als ook voor de hele term \(a+b-c\) te gebruiken.
    Van een kwaliteitsvolle bijlesgever wiskunde mag je verwachten dat hij dergelijke fouten niet maakt.
    
    Dit lijken misschien twee kleine foutjes, maar ze vormen nog maar het topje van de ijsberg.
    Je kunt soms een antwoord krijgen dat iets als \(4ab-4ac\) als ontbonden labelt, puur door je initiële vraag te veranderen.
    Meer zelfs, door in discussie te gaan kun je de chatbot zelfs akkoord laten gaan dat \(x^2-1\) al in ontbonden vorm is!
    Probeer het zelf maar eens bij je favoriete AI-chatbot.
    
    Laat het duidelijk zijn: een chatbot als ChatGPT heeft geen interne definitie van ``ontbonden in factoren'' achter de schermen.
    Ook logische denkregels zijn niet op onomstotelijke wijze geprogrammeerd.
    Het probleem ligt bij hoe een AI als ChatGPT werkt.
 
    Onderliggend wordt een ``Large Language Model'' (LLM) opgesteld.
    Zo'n model kun je vergelijken met een statistisch regressiemodel, waarbij op basis van de onafhankelijke variabele (onze ingegeven tekst, de ``prompt'') een voorspelling wordt gedaan voor de afhankelijke variabele (het bijbehorende tekstuele antwoord).
    
    Dergelijk model wordt opgesteld door informatie te verzamelen in de vorm van een steekproef; bij een LLM is dat een gigantische dataset aan teksten.
    Het gaat ook niet om een soort simpel lineair model met 2 parameters (rico en constante), maar om modellen met meer dan 175 miljard parameters om het meest geschikte antwoord te voorspellen.
    ChatGPT heeft dus geen begrip van (wiskundige) waarheid, accuraatheid, \ldots
    
    Cru geformuleerd: de chatbot verzint telkens het meest geschikte verhaaltje voor mijn input, op basis van de bestudeerde dataset.
    Het mag dus niet verwonderen dat de chatbot fouten maakt: in diens programmatie zijn dat geen fouten maar een antwoord dat meer in lijn ligt met de dataset.
    Zelfs al geef je het model meer parameters en een grotere dataset, het blijft een fundamentele beperking.
    De slaagzekerheid bij dergelijke algebra-vragen zal misschien verhogen (mits een kwalitatieve dataset), maar het zal nooit 100\% zijn.
    In een recente paper (Udandarao, Prabhu et al., 2024) blijkt zelfs dat de performantie slechts logaritmisch stijgt, als de grootte van een LLM lineair toeneemt.
    Misschien komt er zelfs een punt waar het economisch niet interessant is om de accuraatheid van een AI als ChatGPT nog te verhogen.
    
    In feite is een LLM voor dergelijke algebraïsche vragen de ongeschikte tool.
    We zagen eerder dat symbolisch rekenen al decennialang mogelijk is op computers, met een ``Computer Algebra System'' (CAS).
    Zulke systemen kunnen een slaagzekerheid van 100\% halen op algebraïsche rekenvragen; ze passen namelijk letterlijk de algebraïsche rekenregels toe op de input.
    In zulke systemen vind je ook ergens een definitie van ``ontbonden zijn'', en is het onmogelijk om een algebraïsche uitdrukking onterecht als ontbonden te labelen.\\

    Dit betekent echter niet dat er geen mogelijkheden zijn voor AI voor algebra-onderwijs.
    Een LLM is slechts één type AI.
    Je kunt je bijvoorbeeld voorstellen dat ChatGPT beter zou werken met onze opgave, indien de chatbot ook gebruik maakte van een CAS zoals WolframAlpha.
    In dat geval zou ChatGPT onze input vertalen naar een WolframAlpha commando, en zou vervolgens de CAS met 100\% zekerheid zeggen wat de ontbinding van onze uitdrukking is.
    Dergelijke combinaties van verschillende soorten software is exact wat momenteel in de academische wereld onderzocht wordt (Davis, 2024), en wat tegelijk al veel bedrijfjes hopen te kapitaliseren.
    Websites met namen als \href{https://www.google.com/search?q=mathgpt}{``MathGPT''} of ``\href{https://www.google.com/search?q=AI+Math+Tutor}{AI Math Tutor}'' schieten als paddenstoelen uit de grond. 
    Die kunnen vaak het antwoord van een leerling inscannen, wiskundig interpreteren, feedback geven op een specifieke fout en het modelantwoord genereren.
    Dat verloopt dan natuurlijk tegen betaling, wat gemakkelijk goedkoper kan uitkomen dan een bijlesgever betalen, maar dan neemt de leerling er ook het risico op een fout antwoord bij.
    Er gebeuren dus boeiende dingen, zowel in de onderzoeks- als in de bedrijfswereld.
    Misschien komt er een (vrij) betrouwbare AI-chatbot voor wiskunde aan, die een CAS, LLM en nog andere softwaretechnieken combineert, maar wij vinden het vooralsnog te vroeg om hier nog dieper op in te gaan in deze loep.

    Het is in principe nu al mogelijk om AI te gebruiken als ondersteuning voor algebraonderwijs.
    Het brengt echter heel veel dilemma's met zich mee waar de school en vakgroep een antwoord op moet bedenken.
    Enkele eerste hulplijnen hierrond zijn \href{https://www.vlaanderen.be/kenniscentrum-digisprong/themas/innovatie/artificiele-intelligentie}{de visietekst over AI} van het Kenniscentrum Digisprong en \href{https://www.vaia.be/nl/opleidingen}{de website van de Vlaamse AI Academie} die alle nascholingen en contactgegevens van AI-experts bundelt.
    
\end{multicols}

\stopcontents[inhoudloep]		% om inhoud voor het loep artikel 
\vspace{-1cm}
\begin{bronnen}
    \item Bergsen, S. (2023). Automatiseren is meer dan alleen maar oefenen. \textit{Onderwijskennis.nl} (NRO). Geraadpleegd op 11 mei 2024, van \url{www.onderwijskennis.nl/node/4139}
	\item Davis, E. (2024). Mathematics, word problems, common sense, and artificial intelligence. \textit{Bulletin of the American Mathematical Society}. 61 (2024), 287-303, \url{https://doi.org/10.1090/bull/1828}
    \item Deprez, J., \& Op de Beeck, R. (2013). Algebra oefenen met inzicht. \textit{Uitwiskeling}, 29(1), 15-32.  \url{https://www.uitwiskeling.be/2013/01/15/algebra-oefenen-met-inzicht/?v=d3dcf429c679}
    \item Geary, D. C., Boykin, A. W., et al. (2008). Chapter 4: The report of the task group on learning processes, 4-2 to 4-10. U.S. \textit{Department of Education}. \url{https://files.eric.ed.gov/fulltext/ED502980.pdf}
    \item Gevers, P., Leenders, J.H. et al. (1970). Wiskunde 1b. Getallen - meetkunde, \textit{J.B. Wolters Leuven}.
    \item Hartman, J. R., Hart, S. et al. (2023). Designing mathematics standards in agreement with science. \textit{International Electronic Journal of Mathematics Education}, 18(3), em0739.
    \url{https://www.iejme.com/article/designing-mathematics-standards-in-agreement-with-science-13179}
    \item Katholiek Onderwijs Vlaanderen. (2024). Leerplannen. Geraadpleegd op 11 mei 2024 van \url{https://pro.katholiekonderwijs.vlaanderen/ii--wiss-d}
    \item Kilpatrick, J., Swafford, J., \& Findell, B., National Research Council. (2001). Adding It Up: Helping Children Learn Mathematics. Washington, DC: \textit{The National Academies Press}. \url{https://doi.org/10.17226/9822}
    \item KU Leuven. (n.d.). Zelfstudiepakket Wiskunde (Informatica). Geraadpleegd 11 mei 2024, van \url{https://set.kuleuven.be/voorkennis/zomercursus/zomercursusI/rekenen/exercises/begintest_rekenen}
    \item Roelens, M., \& Willems, J. (2008). Algebralessen met applets. \textit{Uitwiskeling}, 24(1), 11-35. \url{https://www.uitwiskeling.be/2008/01/15/algebralessen-met-applets/?v=d3dcf429c679}
    \item Star, J. R., Caronongan, P. et al. (2015). Teaching strategies for improving algebra knowledge in middle and high school students. (NCEE 2014-4333). Washington, DC: \emph{National Center for Education Evaluation and Regional Assistance}. Geraadpleegd op 11 mei 2024, van \url{http://whatworks.ed.gov.}
    \item Tall, D. \& Thomas, M. Encouraging versatile thinking in algebra using the computer. \textit{Educ Stud Math} 22, 125–147 (1991). \url{https://doi.org/10.1007/BF00555720}
    \item Udandarao, V., Prabhu, A. et al. (2024). No ``Zero-Shot'' Without Exponential Data: Pretraining Concept Frequency Determines Multimodal Model Performance. \textit{arXiv preprint arXiv:2404.04125}.
    \item Vanlommel, E., Tytgat, P., De Naeghel, K., Gheysens, L. (2024). Wiskundeplan. Geraadpleegd op 11 mei 2024 van \url{https://wiskundeplan.be/}
\end{bronnen}
\section*{HYPERLINKS}
In de volgorde van het artikel
\begin{itemize}[label={},leftmargin=20pt,itemindent=-20pt,itemsep=6pt,parsep=0pt]%
\item WolframAlpha op \url{https://www.wolframalpha.com}
\item WolfreeAlpha op \url{https://wolfreealpha.onrender.com/}
\item WolfreeAlpha Overzicht Sites op \url{https://www.reddit.com/r/GetStudying/comments/x4luun/get_wolframalpha_pro_feature_for_free/}
\item PhotoMath op \url{https://photomath.com/}
\item WiskundeAcademie op \url{https://www.youtube.com/@WiskundeAcademie}
\item Math with Menno op \url{https://www.youtube.com/@MathwithMenno}
\item KhanAcademy op \url{https://www.khanacademy.org/} 
\item UHasselt cursus op \url{https://sites.google.com/uhasselt.be/platformwiskunde/}
\item UGent cursus op \url{https://elemath.ugent.be/}
\item UAntwerpen cursus op \url{https://www.uantwerpen.be/nl/projecten/aan-de-slag/zelfstudiepakketten/}
\item Algebrakit op \url{https://practice.algebrakit.com/nl-BE/overview}
\item QuickMath op \url{https://quickmath.com}
\item IXL platform op \url{https://uk.ixl.com/maths/algebra}
\item NUMBAS op \url{https://www.numbas.org.uk/}
\item NUMBAS demo op \url{https://numbas.mathcentre.ac.uk/exam/1973/numbas-website-demo/embed/}
\item NUMBAS starthandleiding op \url{https://docs.numbas.org.uk/en/latest/tutorials/write-question.html}
\item NUMBAS originele vraag op \url{https://numbas.mathcentre.ac.uk/question/12000/combining-algebraic-fractions-1/}
\item NUMBAS aangepaste vraag op \url{https://numbas.mathcentre.ac.uk/question/155648/voorbeeld-uitwiskeling/}
\item STACK op \url{https://stack-assessment.org/}
\item ChatGPT op \url{https://chatgpt.com/}
\item Visietekst AI Kenniscentrum Digisprong op \url{https://www.vlaanderen.be/kenniscentrum-digisprong/themas/innovatie/artificiele-intelligentie}
\item Vlaamse AI Academie op \url{https://www.vaia.be/nl/}
\end{itemize}